\chapter*{About this manual}
\addcontentsline{toc}{chapter}{About this manual}

\fbk{} is a zero-barrier workbench for strongly correlated systems and
lanthanide/actinide calculations. It does two things: it \emph{generates}
calculation input files, and it \emph{reads} results (output files, or input
structures) to characterise them. It never runs a calculation for you.

This manual is the complete documentation. It is organised the way the program
is: one chapter per menu entry, each a guided walkthrough. A shorter quick-start
guide (\file{docs/USER\_GUIDE.md}) ships inside the program and is updated with
it; this manual is its expansion.

\section*{How this manual is organised}

\begin{description}[leftmargin=1.6em]
  \item[Part I -- Getting started] what the program is and is not
    (Chapter~\ref{ch:introduction}), how to install it
    (Chapter~\ref{ch:installation}), how to run it and how the
    ``interaction is a script'' model works (Chapter~\ref{ch:running}), and how
    to read the report it writes (Chapter~\ref{ch:report}).
  \item[Part II -- Functions, menu by menu] one chapter per menu number, in
    menu order. Every chapter follows the same seven-part template: purpose;
    what you need; walkthrough; what you get; how to read the numbers; a worked
    example; caveats, refusals and related sections.
  \item[Part III -- Reference] the command line
    (Chapter~\ref{ch:cli}), every criterion and threshold the program uses with
    its provenance (Chapter~\ref{ch:criteria}), and the file formats
    (Chapter~\ref{ch:formats}).
  \item[Part IV -- Development] the architecture, the layering contract, and
    how to extend the rule base, the parsers or the analysers
    (Chapter~\ref{ch:development}).
  \item[Appendices] messages and what to do (Appendix~\ref{app:messages}) and a
    glossary (Appendix~\ref{app:glossary}).
\end{description}

\section*{Conventions}

\begin{itemize}
  \item Menu references: ``\menu{3}'' means the third numbered entry of the
    program's main menu, documented in the chapter of that number. Menu numbers
    are append-only: a number is never reordered or recycled, because it is
    part of the public interface (scripts replay by number).
  \item Session transcripts are typeset in a box and reproduce the program's
    output verbatim, including platform-dependent path separators
    (\code{work\textbackslash file} on Windows, \code{work/file} elsewhere).
  \item Report excerpts are quoted from files the program wrote; the ellipsis
    ``\code{...}'' marks omitted lines.
  \item Every criterion in this manual carries its provenance: a program-manual
    quote (with section and URL), a literature reference (with DOI; the complete
    citation is also printed in the program's reports), or a measured record of
    this project. The three classes are marked \textsf{[Manual]},
    \textsf{[Literature]} and \textsf{[Measured]}.
\end{itemize}

\section*{The examples in this manual are executable}

Every session transcript and report excerpt printed here comes from a real run
of the program. The scripts, the prepared inputs and the captured outputs live
in \file{docs/manual/examples/} inside the repository, and the test suite
replays them: the session captures must reproduce byte for byte
(\file{tests/test\_manual.py}). Nothing in this manual is typed by hand.

\section*{Verification statement}

The verification discipline is the repository's: every parser pattern is
fixed against registered real output (\file{fixtures/}), every rule carries
non-empty evidence, and the acceptance contract is the test suite plus the
layering check (\file{scripts/verify.sh}). Where a statement in this manual
is an inference rather than a documented fact, it says so.
